\chapter{China's Scientific Machine}
\label{c:china-ai4s}

\begin{glance}
If the largest prize is scientific, the question is who is organized for it, and on the primary policy record the answer is uncomfortable: China has been organizing longer, more explicitly and more \emph{industrially} than any Western actor. The State Council's ``AI+'' Opinions place science first of six key actions and chain AI-driven research to ``engineering realisation and product landing,'' which is what makes this industrial policy rather than research funding. The output is measured by the European Commission rather than by Beijing, and it has led the world since 2016. The models are open at 241B and trillion-parameter scale under permissive licences, which is why the case for a Western open scientific-AI consortium cannot be that the commons is missing. Three findings cut the other way and are printed at equal weight: the plan's numeric targets measure diffusion and not discovery; the paradigm-transformation sentence does not survive from the Party's Recommendations into the state Outline; and the newest and most capable Chinese platform has no stated open release at all. \full{25}
\end{glance}

\section{The policy record, read precisely}

The State Council's \emph{Opinions on Deepening the Implementation of the ``AI+'' Initiative}---Guofa [2025] No.~11, signed 21~August 2025---set out six key actions, and ``AI+ science and technology'' is the first of them, ahead of industry~\cite{ai-plus-science} [P]. Ordering in such a document is the sequence in which ministries read their obligations, and putting science ahead of industry in a document whose whole register is industrial says where the leadership believes the compounding lies. The science pillar's third sub-clause is the load-bearing one, and it is why this chapter is titled industrial policy rather than science policy: it directs the integrated development of ``AI-driven technology R\&D, engineering realisation and product landing.'' Discovery is the first stage of a pipeline whose last stage is a product, and the text says so. The lineage runs back to the ministry and science-foundation ``AI for Science'' deployment of March 2023, which Western sources routinely overstate as a plan when it is a deployment of work~\cite{most-ai4s} [P]. The money underneath is bench money spread wide rather than a moonshot: roughly half the project categories in the current science-foundation call foreground AI for science at individual budgets of EUR~240{,}000--600{,}000~\cite{nsfc-ai4s} [A]---enough for students and cluster hours, not frontier training runs: the state supplies compute centrally and grants supply the science.

Two corrections belong in the same section, and both cut against the thesis. First, the numeric commitments in Guofa [2025] No.~11 are \emph{diffusion} targets---intelligent-terminal and agent penetration above 70\% by 2027 and 90\% by 2030---and there is \textbf{no numeric target for science at all}~\cite{ai-plus-science} [P]. A planning system whose distinctive competence is hitting what it measures has, in its flagship AI document, measured adoption and declined to measure discovery. Second, and more precisely: the Party's \emph{Recommendations} for the 15th Five-Year Plan, released 28~October 2025, contain the instruction to \emph{use artificial intelligence to lead a transformation of the scientific research paradigm}---and \textbf{that phrase does not survive into the Outline adopted by the National People's Congress in March 2026}, where AI appears under compute, algorithm and data supply while basic research appears in its own chapter without AI attached~\cite{fyp15-ai4s} [P]. The Recommendations state ambition; the Outline allocates. The deflationary reading is that a paradigm shift is not a line item; the consequential reading is that the drafting converted \emph{how research should be done} into \emph{what infrastructure should be bought}, and the state's own prize list agrees with it. Either way the gap between Party ambition and state programming is real, documentable, and the strongest single piece of evidence against this chapter.

\section{The output, the artifacts, and the machine that won the prize}

The most useful measurement was made in Brussels, which is the reason to use it. The European Commission's bibliometric study of some three million papers across eighty fields finds that \textbf{China overtook both the United States and the European Union in AI-assisted scientific output from 2016}, producing more than 25{,}000 AI-assisted papers in 2022 against roughly 15{,}000 for the EU and 12{,}000 for the US, with reported advantages on novelty and citation impact as well as volume~\cite{ec-ai4s-biblio} [V on the counts, A on the impact findings]; the Nature Index points the same way~\cite{nature-index-2026}. Two limits belong beside it. \emph{Output is not discovery}---papers are the first term of the funnel of Chapter~\ref{c:science}, whose later terms are all fractions below one. And \emph{the window ends in 2022}, before the model generation this book is about, so the study measures the run-up rather than the present state: the position was built before the current phase, not during it.

The artifacts matter more, and the fact that matters most about them is not any benchmark number but their licences. Shanghai AI Laboratory's \textbf{Intern-S1} carries 241B parameters pre-trained on 5T tokens of which \emph{2.5T are scientific}---a resource-allocation decision no Western laboratory has publicly matched---and is published under \textbf{Apache~2.0}; its scores are self-reported and recorded at that grade~\cite{intern-s1-2025} [P]. ByteDance's \textbf{Protenix} is released under the same licence with code, weights and pipelines, claiming to be the first fully open-source model to outperform AlphaFold3 at matched data cutoff, scale and inference budget~\cite{protenix} [P on the release, developer-evaluated on the comparison]---and the significance is the licence, not the benchmark. The open-weight ratchet of Chapter~\ref{c:elements} has reached structural biology, and it reached it from a Chinese laboratory; a database is a result and a weight file is a means of production, and the West's most diffused scientific asset was given away as the former. \textbf{DeepModeling and DP Technology} are the deepest instance: DeePMD-kit, DPA-2 covering 73 elements across eighteen pre-training datasets, the ABACUS electronic-structure package, the OpenLAM initiative, more than ten thousand community members, with code, weights and datasets all open~\cite{deepmodeling} [P]. On the record at \datafreeze\ that is among the most genuinely open scientific-AI ecosystems anywhere, Western or Chinese---a sentence that cuts against the framing a Western reader will bring. Against which stands the counter-observation, made prominently because it must be: the newest and most capable platform---Intern Duanyan, announced July 2026, a five-layer discovery stack with a dry-lab/wet-lab closed loop---has \textbf{no stated open-source release} and is access-by-invitation~\cite{intern-s1-2025} [P]. The ratchet compels weights onto the commons because adoption is the currency; it compels nothing whatever about a platform, whose rent lives in orchestration and instrument access rather than in parameters. Openness here may be a phase rather than a principle. The same ambiguity runs through the institutions: CAS ScienceOne is built on \emph{Chinese} open foundation models and consumes AlphaFold and MatterGen as tools underneath its own orchestration layer---the mirror mechanism of Chapter~\ref{c:mirror} in a scientific setting---while its parameters, training compute and open-source status are undisclosed and its performance claims self-reported, so this book credits its architecture and not its capability~\cite{cas-scienceone} [P].

The infrastructure supplies the tell. Pengcheng Cloudbrain, the first fully domestic large-scale intelligent computing system, was announced in July 2026 as winner of the \textbf{National Science and Technology Progress Award, First Prize}, and is reported to support training for more than half of China's domestic foundation models~\cite{pengcheng-cloudbrain} [P on the award, A on the ratio]. \emph{The state's top science prize went to a computer, not to a discovery made on one.} A national prize list is a revealed preference over what a state believes is scarce, and this one says the scarce thing is not a result but the domestic capacity to produce results---the argument of Chapter~\ref{c:elements} restated by the awarding body. Western prize lists decorate the structure, the forecast, the molecule, and treat the machine as procurement; neither choice is obviously right, but the asymmetry is real.

\section{Scored on the playbook, and what it does to the consortium}

On the playbook of Chapter~\ref{c:precedents} the programme reads cleanly: target selection aimed at the research paradigm itself; demand before supply, with the Outline naming government \emph{purchase and leasing} of compute and universal accessibility for small users; selection rather than planning, in six-plus institutions with overlapping mandates and no designated national laboratory; supply-chain capture running model, framework, code, chip and machine, with a domestic electronic-structure code inserted into the label pipeline of the flagship domestic atomic model~\cite{fyp15-ai4s,deepmodeling,ai4s-industrial}. Two steps are not scored: overcapacity has the shape but no published utilization figures, so overbuild cannot be told from stranded provincial assets [S], and gatekeeping is not observed at all---no restriction on scientific model release, no export control on scientific weights. The industrial hinge is the distinctive finding: of ten flagship 2026 results, three are aerospace design, industrial-chemistry formulation and an analogue chip-design agent, the flywheel aimed directly at the layer Chapter~\ref{c:elements} identifies as the binding physical constraint---while three of the same ten carry Western co-authors, so any account treating Chinese AI for science as a hermetically closed national system is wrong on the byline evidence~\cite{ai4s-industrial} [P].

Which sharpens the consortium question of Chapter~\ref{c:science} rather than settling it. The case usually offered for a Western or neutral open scientific-AI consortium---that the commons is thin and public money must create an open base---is, on this evidence, false: the open base exists, and the researchers a consortium proposes to serve are already downloading it. The defensible case is narrower and stronger: \emph{the commons has one principal contributor and no governance the rest of the world is party to}---contribution and governance being different goods. Each of the three insurance risks is sharpened here: continuity, by an unopened platform at the top of the stack; verification, because four of the most consequential artifacts in the field cannot be checked in the way Chapter~\ref{c:trust} requires; monoculture, if one national ecosystem supplies the potentials, the scientific foundation models and the agent frameworks a generation of results depends on. The opposing case is strong: a consortium built to \emph{exclude} its largest contributor is a worse instrument than one built to include it, forfeiting both the artifacts and the neutrality that is its only durable asset; and if scientific knowledge is non-rival, a policy that slows the commons in order to change its ownership destroys more value than it redistributes, in the one domain where the social return runs above four dollars on the dollar. This chapter does not resolve that; it establishes the question: whether an open commons with one principal contributor and no shared governance is a commons or a dependency, and what instrument short of exclusion changes the answer.
